\documentclass[11pt]{amsart}

\usepackage[T1]{fontenc}
\usepackage[utf8]{inputenc}
% See the note in the companion article: without an outline font the T1 text fonts are
% embedded as Type 3 bitmaps and the ligatures cease to be searchable in the PDF.
\usepackage{lmodern}
\input{glyphtounicode}
\pdfgentounicode=1

\usepackage{amsmath,amssymb,amsthm,mathtools}
\usepackage{array,booktabs,tabularx}
\usepackage{xcolor}
\usepackage{geometry}
\usepackage{enumitem}
\usepackage{hyperref}

\geometry{margin=1in}
\setlength{\emergencystretch}{3em}
\allowdisplaybreaks
\hypersetup{
  colorlinks=true,
  linkcolor=blue!45!black,
  citecolor=blue!45!black,
  urlcolor=blue!45!black,
  pdftitle={Kernel-Defined Multilinear Operators and Conditional Analytic Extensions of a Finite Projection Theory},
  pdfauthor={Eliuth Chavero Jasso},
  pdfsubject={Multilinear integral operators on L2, associator kernels, truncated Hodge compatibility, induced Markov operators},
  pdfkeywords={multilinear integral operator, associator, orthogonal projection, cochain complex, Hodge Laplacian, spectral truncation, Markov operator}
}

\newtheorem{theorem}{Theorem}[section]
\newtheorem{proposition}[theorem]{Proposition}
\newtheorem{corollary}[theorem]{Corollary}
\newtheorem{lemma}[theorem]{Lemma}
\theoremstyle{definition}
\newtheorem{definition}[theorem]{Definition}
\newtheorem{construction}[theorem]{Construction}
\newtheorem{example}[theorem]{Example}
\theoremstyle{remark}
\newtheorem{remark}[theorem]{Remark}
\newtheorem{question}[theorem]{Question}

\newcommand{\ran}{\operatorname{ran}}
\newcommand{\Tr}{\operatorname{Tr}}
\newcommand{\op}{\mathrm{op}}
\newcommand{\norm}[1]{\lVert#1\rVert}
\newcommand{\Types}{\mathcal T}
\newcommand{\Kop}{\mathcal K}
\newcommand{\Energy}{\mathfrak E}
\newcommand{\proj}{\mathrm{proj}}
\newcommand{\amb}{\mathrm{amb}}
\newcommand{\red}{\mathrm{red}}

\title[Kernel-defined multilinear operators]{Kernel-Defined Multilinear Operators\\ and Conditional Analytic Extensions\\ of a Finite Projection Theory}
\author{Eliuth Chavero Jasso}
\address{Independent Researcher, Apizaco, Tlaxcala, Mexico}
\thanks{Companion to \cite{ChaveroJassoTrees}, which contains the finite-dimensional
theory. Nothing in the present article has been checked by anyone other than the author,
and no claim of originality is made.}
\date{31 July 2026}

\begin{document}
\maketitle

\begin{abstract}
A companion article \cite{ChaveroJassoTrees} studies finite composition trees of
multilinear maps between finite-dimensional Hilbert spaces whose intermediate outputs are
orthogonally projected onto prescribed subspaces, and proves an exact local error
decomposition together with uniform coefficients $k$ and $k-1$ for the ambient and
projected root errors. The present article asks which parts of that setting survive when
the finite-dimensional spaces are replaced by $L^2$ spaces and the multilinear maps by
integral operators, and it is explicit at every step about what is proved, what is proved
only under stated hypotheses, and what is merely proposed.

We prove that a kernel in $L^2(X^{a+1})$ defines a bounded multilinear operator on
$L^2(X)$, and --- this closes a gap in an earlier formulation --- that the composite
kernels obtained by inserting one such operator into another again lie in $L^2$, with
$\norm{\kappa_L}_{L^2(X^{6})}\le\norm{\kappa}_{L^2(X^{4})}^{2}$ in the ternary case, so
that the associator of a kernel-defined operator is itself given by a square-integrable
defect kernel. We prove that the associator characterises this defect: under separability
of $L^2(X,\nu)$, the associator vanishes on all products of $L^2$ functions if and only if
the defect kernel vanishes almost everywhere. We record the exact identity relating an
associator to the commutator defect of left multiplication operators, and we are careful
to present the resulting algebraic tensor as a definition rather than as a curvature.
For finite cochain complexes we give the standard conditions under which a degree-zero
graded operator descends to cohomology and commutes with a Hodge Laplacian, and we state
the spectral-truncation version as a statement about the truncation only. An induced
Markov operator and its Dirichlet form are presented as a conditional construction whose
hypothesis we do not verify for any kernel. Continuum limits, membership in a
pseudodifferential class, and microlocal regularity are open questions and are collected
as such.
\end{abstract}

\subjclass[2020]{47H60, 47G10, 46E30, 58J50, 17A30}
\keywords{multilinear integral operator, associator kernel, orthogonal projection, cochain
complex, Hodge Laplacian, spectral truncation, Markov operator}

\tableofcontents

\section{Introduction}
\label{sec:intro}

\subsection{What the companion article establishes}
\label{sec:finite}

Let $\Types$ be a finite set of types, and for each $\tau\in\Types$ let $V_\tau$ be a
finite-dimensional Hilbert space, $Q_\tau:W_\tau\to V_\tau$ an isometry and
$P_\tau=Q_\tau Q_\tau^{*}$ the associated orthogonal projector. A finite rooted ordered
tree whose internal vertices carry multilinear maps
$\mu_v:\prod_jV_{\tau(v,j)}\to V_{\tau(v)}$ admits two evaluations: the unprojected one,
and the one in which the output of every internal vertex is replaced by its orthogonal
projection. Writing $E_T^{\amb}$, $E_T^{\proj}$, $E_T^{\perp}$ and $E_T^{\red}$ for the
four resulting measures of discrepancy at the root, the companion article proves that
\[
 \bigl(E_T^{\amb}\bigr)^2=\bigl(E_T^{\proj}\bigr)^2+\bigl(E_T^{\perp}\bigr)^2,
 \qquad
 E_T^{\red}=E_T^{\proj},
\]
that the discrepancy at each vertex admits an exact decomposition into a locally created
residual and a term for every nonempty subset of the children carrying an error, and that
under uniform bounds $\norm{\mu_v}_\op\le M$ and $\norm{r_v}_\op\le\rho$ a tree with
$k\ge1$ internal vertices satisfies
\begin{equation}
\label{eq:companion}
 E_T^{\amb}\le k\rho M^{k-1}L_T,
 \qquad
 E_T^{\proj}=E_T^{\red}\le(k-1)\rho M^{k-1}L_T,
\end{equation}
$L_T$ being the product of the leaf norms. The coefficient $k-1$ is an upper bound; the
companion article does not show it is attained at fixed $\rho/M>0$, and leaves the exact
constant open already for two internal vertices.

\emph{We reproduce none of these proofs.} They are stated above only to fix the notation
and the hypotheses that the present article refers to, and every use of them below is a
citation of \cite{ChaveroJassoTrees}.
\label{sec:finite-end}

\subsection{What this article does}

The finite theory has three features that are not obviously stable under passage to
infinite dimensions: it uses operator norms of multilinear maps, it uses orthogonal
projectors with closed range, and it uses composition of multilinear maps. Each raises a
separate analytic question, and the answers are of different quality.

Section~\ref{sec:kernel} shows that multilinear operators defined by square-integrable
kernels behave well: they are bounded, they depend continuously on the kernel, and ---
Lemma~\ref{lem:composite} --- their composites are again defined by square-integrable
kernels. This last point matters because the earlier formulation of this material defined
the associator kernel without establishing that it exists, and the estimate supplied here
removes that gap.

Section~\ref{sec:associator} treats the associator. Proposition~\ref{prop:converse} shows
that the associator determines the defect kernel exactly, under a separability hypothesis
that we state explicitly because it is not automatic. Section~\ref{sec:curvature} records
the identity relating the associator to the commutator defect of left multiplication
operators. We emphasise that the algebraic tensor obtained from an associator is a
definition, not a curvature tensor of a connection.

Sections~\ref{sec:variational} and~\ref{sec:cohomology} concern structures that are
finite or truncated: a variational formulation on a Stiefel manifold, and the descent of a
degree-zero graded operator to the cohomology of a finite cochain complex, together with
its spectral-truncation analogue. These are standard statements; we give them because the
truncated versions are what the computational work of the companion program actually
tests, and it is important that they not be mistaken for continuum statements.

Section~\ref{sec:markov} presents an induced Markov operator and its Dirichlet form as a
\emph{conditional construction}. Its hypothesis --- that a contraction of the kernel
produces a symmetric nonnegative weight --- is not verified here for any kernel, and we
say so.

Section~\ref{sec:open} collects the open analytical questions. Continuum limits of the
finite theory, membership of kernel-defined operators in a multilinear pseudodifferential
class, and microlocal regularity are questions, not results, and no claim about them is
made anywhere in this article.

\section{Multilinear operators defined by square-integrable kernels}
\label{sec:kernel}

Let $(X,\nu)$ be a $\sigma$-finite measure space and $H=L^2(X,\nu)$. Throughout, $\nu$ is
the measure; the letter $\mu$ is reserved for the multilinear maps of the finite theory
and does not appear in this section.

\begin{definition}[Kernel-defined multilinear operator]
\label{def:kernelop}
For $\kappa\in L^{2}(X^{a+1},\nu^{\otimes(a+1)})$ define
\[
 \Kop_\kappa(f_1,\ldots,f_a)(p)
 =\int_{X^{a}}\kappa(p;q_1,\ldots,q_a)\prod_{j=1}^{a}f_j(q_j)\,
 d\nu^{\otimes a}(q_1,\ldots,q_a).
\]
\end{definition}

\begin{proposition}[Boundedness]
\label{prop:bounded}
If $\kappa\in L^{2}(X^{a+1})$ then $\Kop_\kappa$ is a well-defined bounded $a$-linear map
$H^{a}\to H$, and
\[
 \norm{\Kop_\kappa(f_1,\ldots,f_a)}_{2}
 \le\norm{\kappa}_{L^{2}(X^{a+1})}\prod_{j=1}^{a}\norm{f_j}_{2}.
\]
\end{proposition}

\begin{proof}
By Tonelli's theorem, $p\mapsto\norm{\kappa(p;\cdot)}_{L^{2}(X^{a})}$ is measurable and
\[
 \int_X\norm{\kappa(p;\cdot)}_{L^{2}(X^{a})}^{2}\,d\nu(p)=\norm{\kappa}_{L^{2}(X^{a+1})}^{2}.
\]
For $\nu$-almost every $p$ the function $\kappa(p;\cdot)$ lies in $L^{2}(X^{a})$, and
$f_1\otimes\cdots\otimes f_a\in L^{2}(X^{a})$ with norm $\prod_j\norm{f_j}_2$ by Tonelli
again. Cauchy--Schwarz on $X^{a}$ gives, for almost every $p$, both absolute convergence
of the defining integral and
\[
 \bigl\lvert\Kop_\kappa(f_1,\ldots,f_a)(p)\bigr\rvert
 \le\norm{\kappa(p;\cdot)}_{L^{2}(X^{a})}\prod_{j}\norm{f_j}_{2}.
\]
Squaring and integrating in $p$ gives the stated estimate.
\end{proof}

\begin{proposition}[Dependence on the kernel]
\label{prop:stability}
For $\kappa,\widetilde\kappa\in L^{2}(X^{a+1})$,
\[
 \norm{\Kop_\kappa(f_1,\ldots,f_a)-\Kop_{\widetilde\kappa}(f_1,\ldots,f_a)}_{2}
 \le\norm{\kappa-\widetilde\kappa}_{L^{2}(X^{a+1})}\prod_j\norm{f_j}_2 .
\]
In particular $\kappa_n\to\kappa$ in $L^{2}(X^{a+1})$ implies
$\Kop_{\kappa_n}\to\Kop_\kappa$ in the multilinear operator norm.
\end{proposition}

\begin{proof}
$\kappa\mapsto\Kop_\kappa$ is linear, so this is Proposition~\ref{prop:bounded} applied to
$\kappa-\widetilde\kappa$. The second assertion follows because the estimate is uniform
over unit vectors.
\end{proof}

\subsection{Composite kernels}

We now compose two ternary operators; the general $a$-ary case is identical.

\begin{lemma}[Composite kernels are square-integrable]
\label{lem:composite}
Let $\kappa\in L^{2}(X^{4})$ and define, for almost every argument,
\[
 \kappa_L(p;q_1,\ldots,q_5)
 =\int_X\kappa(p;u,q_4,q_5)\,\kappa(u;q_1,q_2,q_3)\,d\nu(u),
\]
\[
 \kappa_R(p;q_1,\ldots,q_5)
 =\int_X\kappa(p;q_1,q_2,u)\,\kappa(u;q_3,q_4,q_5)\,d\nu(u).
\]
Then the defining integrals converge absolutely for
$\nu^{\otimes6}$-almost every $(p;q_1,\ldots,q_5)$, both $\kappa_L$ and $\kappa_R$ belong
to $L^{2}(X^{6})$, and
\[
 \norm{\kappa_L}_{L^{2}(X^{6})}\le\norm{\kappa}_{L^{2}(X^{4})}^{2},
 \qquad
 \norm{\kappa_R}_{L^{2}(X^{6})}\le\norm{\kappa}_{L^{2}(X^{4})}^{2}.
\]
\end{lemma}

\begin{proof}
Write $g(p,q_4,q_5)=\norm{\kappa(p;\cdot,q_4,q_5)}_{L^{2}(X)}$ and
$h(q_1,q_2,q_3)=\norm{\kappa(\cdot;q_1,q_2,q_3)}_{L^{2}(X)}$, both measurable by Tonelli.
Cauchy--Schwarz in the variable $u$ gives, pointwise,
\[
 \int_X\lvert\kappa(p;u,q_4,q_5)\kappa(u;q_1,q_2,q_3)\rvert\,d\nu(u)
 \le g(p,q_4,q_5)\,h(q_1,q_2,q_3),
\]
so the integral defining $\kappa_L$ converges absolutely whenever the right-hand side is
finite, and
\[
 \lvert\kappa_L(p;q_1,\ldots,q_5)\rvert^{2}\le g(p,q_4,q_5)^{2}\,h(q_1,q_2,q_3)^{2}.
\]
Integrating over $X^{6}$ and using Tonelli to separate the variables,
\[
 \norm{\kappa_L}_{L^{2}(X^{6})}^{2}
 \le\Bigl(\int_{X^{3}}g^{2}\,d\nu^{\otimes3}\Bigr)
 \Bigl(\int_{X^{3}}h^{2}\,d\nu^{\otimes3}\Bigr)
 =\norm{\kappa}_{L^{2}(X^{4})}^{2}\cdot\norm{\kappa}_{L^{2}(X^{4})}^{2},
\]
since each factor is $\norm{\kappa}_{L^{2}(X^{4})}^{2}$ by Tonelli. Finiteness of the
left-hand side forces the right-hand side of the pointwise estimate to be finite almost
everywhere, which is the asserted almost-everywhere absolute convergence. The argument for
$\kappa_R$ is the same with the roles of the variables exchanged.
\end{proof}

\begin{remark}
\label{rem:gapclosed}
Lemma~\ref{lem:composite} is what makes Definition~\ref{def:assockernel} below legitimate.
An earlier formulation of this material introduced $\kappa_L$, $\kappa_R$ and their
difference, and defined an energy as the squared $L^{2}(X^{6})$ norm of that difference,
without establishing that any of the three exists. Square-integrability of a kernel does
not in general pass to composites of the associated operators for free; here it does, and
the reason is the elementary estimate above.
\end{remark}

\section{The associator of a kernel-defined operator}
\label{sec:associator}

For a ternary operation there is no canonical associator without a choice of how the
compositions are inserted. We fix one, and use it consistently.

\begin{definition}[Five-input ternary associator]
\label{def:assoc5}
For a ternary multilinear map $\mu$ on a vector space,
\[
 A^{(5)}_{\mu}(x_1,\ldots,x_5)
 =\mu(\mu(x_1,x_2,x_3),x_4,x_5)-\mu(x_1,x_2,\mu(x_3,x_4,x_5)).
\]
In operadic notation, $A^{(5)}_\mu=\mu\circ_1\mu-\mu\circ_3\mu$ \cite{Yau2016,LodayVallette2012}.
\end{definition}

Other conventions --- the binary associator induced by anchoring one argument, or
insertion at another slot --- have different domains and different meanings. They must not
be merged into a single quantity, and we do not do so.

\begin{definition}[Defect kernel and associator energy]
\label{def:assockernel}
For $\kappa\in L^{2}(X^{4})$ put $\Phi_\kappa=\kappa_L-\kappa_R$, which lies in
$L^{2}(X^{6})$ by Lemma~\ref{lem:composite}, with
$\norm{\Phi_\kappa}_{L^{2}(X^{6})}\le2\norm{\kappa}_{L^{2}(X^{4})}^{2}$. The
\emph{associator energy} is
\[
 \Energy_A(\kappa)=\norm{\Phi_\kappa}_{L^{2}(X^{6})}^{2}
 \;\le\;4\norm{\kappa}_{L^{2}(X^{4})}^{4}.
\]
\end{definition}

\begin{proposition}[The defect kernel represents the associator]
\label{prop:represents}
For all $f_1,\ldots,f_5\in L^{2}(X)$,
\[
 A^{(5)}_{\Kop_\kappa}(f_1,\ldots,f_5)(p)
 =\int_{X^{5}}\Phi_\kappa(p;q_1,\ldots,q_5)\prod_{j=1}^{5}f_j(q_j)\,d\nu^{\otimes5}(q)
\]
for $\nu$-almost every $p$; that is, $A^{(5)}_{\Kop_\kappa}=\Kop_{\Phi_\kappa}$.
\end{proposition}

\begin{proof}
By Lemma~\ref{lem:composite} the function
$(u,q)\mapsto\kappa(p;u,q_4,q_5)\kappa(u;q_1,q_2,q_3)\prod_jf_j(q_j)$ is absolutely
integrable over $X^{6}$ for almost every $p$: indeed Cauchy--Schwarz bounds its integral
by $\norm{\kappa(p;\cdot)}_{L^{2}(X^{3})}\norm{\kappa}_{L^{2}(X^{4})}\prod_j\norm{f_j}_2$,
which is finite for almost every $p$. Fubini's theorem therefore permits the interchange
of the $u$-integration with the $q$-integration, which converts
$\Kop_\kappa(\Kop_\kappa(f_1,f_2,f_3),f_4,f_5)$ into integration against $\kappa_L$. The
same argument converts $\Kop_\kappa(f_1,f_2,\Kop_\kappa(f_3,f_4,f_5))$ into integration
against $\kappa_R$. Subtracting gives the claim.
\end{proof}

\begin{proposition}[The associator determines the defect kernel]
\label{prop:converse}
Assume in addition that $L^{2}(X,\nu)$ is separable. Then
\[
 A^{(5)}_{\Kop_\kappa}(f_1,\ldots,f_5)=0
 \text{ in }L^{2}(X)\text{ for all }f_1,\ldots,f_5\in L^{2}(X)
 \iff
 \Phi_\kappa=0
 \quad\nu^{\otimes6}\text{-a.e.}
\]
\end{proposition}

\begin{proof}
If $\Phi_\kappa=0$ almost everywhere then the associator vanishes by
Proposition~\ref{prop:represents}.

Conversely, let $\{g_m\}_{m\in\mathbb N}$ be a countable dense subset of $L^{2}(X,\nu)$,
which exists by separability. For each of the countably many $5$-tuples
$(m_1,\ldots,m_5)$, Proposition~\ref{prop:represents} and the hypothesis give
\[
 \int_{X^{5}}\Phi_\kappa(p;q)\prod_{j=1}^{5}g_{m_j}(q_j)\,d\nu^{\otimes5}(q)=0
\]
for $p$ outside a $\nu$-null set $N_{m_1\cdots m_5}$. Let
$N=\bigcup N_{m_1\cdots m_5}$, a countable union of null sets, hence null. Enlarging $N$
if necessary, we may also assume that $\Phi_\kappa(p;\cdot)\in L^{2}(X^{5})$ for
$p\notin N$, which holds almost everywhere by Lemma~\ref{lem:composite} and Tonelli.

Fix $p\notin N$. The linear span of $\{g_{m_1}\otimes\cdots\otimes g_{m_5}\}$ is dense in
$L^{2}(X^{5})$: finite linear combinations of products of elements of a dense subset of
$L^{2}(X)$ are dense in the Hilbert space tensor product $L^{2}(X)^{\otimes5}\cong
L^{2}(X^{5})$ \cite[Sec.~II.4]{ReedSimon1980}. The function $\Phi_\kappa(p;\cdot)$ is an
element of $L^{2}(X^{5})$ orthogonal to a dense subspace, hence zero. As this holds for
every $p\notin N$, we get $\Phi_\kappa=0$ almost everywhere by Tonelli.
\end{proof}

\begin{remark}
\label{rem:separability}
Separability is a genuine hypothesis and is not automatic for a $\sigma$-finite measure
space; it holds when the underlying $\sigma$-algebra is countably generated modulo null
sets. Without it the countable-dense-family argument is unavailable, and we make no claim.
\end{remark}

\section{Associators and commutator defects}
\label{sec:curvature}

Let $\circ$ be any bilinear operation on a vector space $V$ --- no algebra axioms are
assumed. Write $\mathsf L_xz=x\circ z$ for left multiplication, $[x,y]=x\circ y-y\circ x$,
and
\[
 A(x,y,z)=(x\circ y)\circ z-x\circ(y\circ z),
 \qquad
 R_{\mathrm{std}}(x,y)=[\mathsf L_x,\mathsf L_y]-\mathsf L_{[x,y]}.
\]

\begin{theorem}[Commutator defect and associator]
\label{thm:commutator}
For all $x,y,z\in V$,
\[
 R_{\mathrm{std}}(x,y)z=A(y,x,z)-A(x,y,z).
\]
\end{theorem}

\begin{proof}
Expanding the definitions,
\begin{align*}
 R_{\mathrm{std}}(x,y)z
 &=x\circ(y\circ z)-y\circ(x\circ z)-(x\circ y)\circ z+(y\circ x)\circ z\\
 &=\bigl[x\circ(y\circ z)-(x\circ y)\circ z\bigr]
 -\bigl[y\circ(x\circ z)-(y\circ x)\circ z\bigr]\\
 &=-A(x,y,z)+A(y,x,z).\qedhere
\end{align*}
\end{proof}

\begin{definition}[Associator-based algebraic tensor]
\label{def:ralg}
One may \emph{define}, by convention, $R_{\mathrm{alg}}:=A$.
\end{definition}

\begin{remark}[This is a definition, not a curvature theorem]
\label{rem:notcurvature}
Definition~\ref{def:ralg} assigns a name to the associator. It does not identify
$R_{\mathrm{alg}}$ with the curvature tensor of a connection: no connection, no bundle and
no covariant derivative have been introduced, and none is implied. The only universal
identity available is the antisymmetrised relation of Theorem~\ref{thm:commutator}. An
outright identity $R_{\mathrm{std}}=A$ would require additional hypotheses, particular
symmetries, or a different convention, and we do not assert one.
\end{remark}

\begin{remark}[Distinct identities]
\label{rem:identities}
Associativity, cyclic symmetry, the Jacobi identity and Filippov's fundamental identity
\cite{Filippov1985} are distinct conditions. A small residual for one of them implies
nothing about the others, and residuals for different conventions must not be aggregated
into a single score. The companion article's numerical study of signed combinations
\cite{ChaveroJassoTrees} keeps them separate for this reason.
\end{remark}

\section{Reduced operators, closure leakage, and a variational formulation}
\label{sec:variational}

\subsection{Reduced operators}

Let $W$ be a Hilbert space, $Q:W\to V$ an isometry and $P=QQ^{*}$. For a multilinear map
$\mu$ on $V$ the \emph{reduced map} is
$\bar\mu(z_1,\ldots,z_a)=Q^{*}\mu(Qz_1,\ldots,Qz_a)$, with ambient realisation
$Q\bar\mu(z_1,\ldots,z_a)=P\mu(Qz_1,\ldots,Qz_a)$. The subspace family is invariant under
$\mu$ exactly when the \emph{closure-leakage map}
$\mathcal C_\mu=(I-P)\mu(P\,\cdot,\ldots,P\,\cdot)$ vanishes identically; its operator
norm $\rho_\mu=\norm{\mathcal C_\mu}_{\op}$ is the quantity denoted $\rho$ in
\eqref{eq:companion}.

\begin{remark}[An operator norm is not a sampled average]
\label{rem:sampled}
Numerical work frequently reports a normalised sampled quantity such as
\[
 \Energy_{\mathrm{closure}}
 =\mathbb E\left[
 \frac{\norm{(I-P)\mu(PX_1,\ldots,PX_a)}^{2}}
 {\norm{\mu(PX_1,\ldots,PX_a)}^{2}+\varepsilon}\right]
\]
for random $X_i$. This is not $\rho_\mu$, and a small value of it does not bound
$\rho_\mu$. The two must be reported separately.
\end{remark}

\subsection{Low-rank representation}

To avoid storing a coefficient tensor of size $d_0d_1\cdots d_a$ one may use a rank-$R$
canonical polyadic representation \cite{KoldaBader2009}. Such a representation is not
unique: factors may be rescaled subject to a product constraint, and components permuted,
without changing the tensor. Any comparison of factors must therefore be carried out
modulo scaling, phase and permutation, and identifiability requires hypotheses of Kruskal
type \cite{Kruskal1977}. A low-rank representation is a computational device; on its own
it introduces no algebraic structure.

\subsection{Variational formulation on a Stiefel manifold}

Parametrising $P=QQ^{*}$ with $Q^{*}Q=I_r$ places $Q$ on the Stiefel manifold
\[
 \mathrm{St}(d,r)=\{Q\in\mathbb K^{d\times r}:Q^{*}Q=I_r\},
 \qquad
 T_Q\mathrm{St}(d,r)=\{Z:Q^{*}Z+Z^{*}Q=0\}.
\]
With respect to the metric induced by the Euclidean embedding, the Riemannian gradient of
a smooth objective $\Energy$ with Euclidean gradient $G$ is
\[
 \operatorname{grad}\Energy(Q)=G-Q\operatorname{sym}(Q^{*}G),
 \qquad
 \operatorname{sym}(B)=\tfrac12(B+B^{*}),
\]
see \cite[Sec.~2.4]{EdelmanAriasSmith1998} and \cite[Ch.~3]{AbsilMahonySepulchre2008}. The
canonical metric on $\mathrm{St}(d,r)$ gives a different gradient; we state the embedded
Euclidean one because that is the one used in the accompanying computations.

\begin{remark}
\label{rem:localmin}
A numerically located stationary point of such an objective need not be a global minimum,
and the subspace it determines need not be canonical in any sense. This is the ordinary
situation for nonconvex optimisation on a manifold, and no result below depends on the
contrary.
\end{remark}

The tangent space $T_Q\mathrm{St}(d,r)$ is the one place in this article where the word
\emph{tangent} refers to a tangent space. Elsewhere, $\ran P$ is a fixed linear subspace
and the words \emph{projected} and \emph{orthogonal} are used.

\section{Finite cochain complexes and spectral truncation}
\label{sec:cohomology}

\subsection{Descent to cohomology}

Let $C^{\bullet}:C^{0}\xrightarrow{d_0}C^{1}\xrightarrow{d_1}\cdots\to C^{n}$ be a finite
complex of finite-dimensional inner-product spaces with $d_{p+1}d_p=0$, and
$H^{p}(C,d)=\ker d_p/\operatorname{im}d_{p-1}$.

\begin{proposition}[Descent]
\label{prop:descent}
Let $S=(S_p)_{p}$ be a graded operator \emph{of degree zero}, i.e.\ $S_p:C^{p}\to C^{p}$
for each $p$. If $S_{p+1}d_p=d_pS_p$ for every $p$, then $S$ induces a well-defined map
$S_{*}:H^{p}(C,d)\to H^{p}(C,d)$ for every $p$.
\end{proposition}

\begin{proof}
If $x\in\ker d_p$ then $d_pS_px=S_{p+1}d_px=0$, so $S_p$ preserves cocycles. If
$x=d_{p-1}y$ then $S_px=S_pd_{p-1}y=d_{p-1}S_{p-1}y$, so $S_p$ preserves coboundaries.
Hence $[x]\mapsto[S_px]$ is well defined.
\end{proof}

The degree-zero hypothesis is required and is stated explicitly; without it the two
displayed compatibilities do not even typecheck.

\subsection{Hodge compatibility}

With $d^{*}$ the adjoints and $\Delta_p=d_{p-1}d_{p-1}^{*}+d_p^{*}d_p$ the Hodge
Laplacian, put $\mathcal H^{p}=\ker\Delta_p$.

\begin{proposition}[Harmonic subspaces]
\label{prop:hodge}
If $S$ commutes with $d$ and with $d^{*}$ in the graded sense, then $S$ commutes with
$\Delta$, hence $S_p(\mathcal H^{p})\subseteq\mathcal H^{p}$.
\end{proposition}

\begin{proof}
Immediate from $\Delta=dd^{*}+d^{*}d$ and bilinearity of the commutator. The last
assertion follows because $\Delta_pS_px=S_p\Delta_px=0$ for $x\in\mathcal H^{p}$.
\end{proof}

For a finite complex, $\mathcal H^{p}\cong H^{p}(C,d)$ by the finite-dimensional Hodge
decomposition; on a compact Riemannian manifold the corresponding statement for the de
Rham complex is the Hodge theorem \cite[Ch.~6]{Warner1983}. The two settings are different
and we do not pass from one to the other.

\subsection{Spectral truncation}

Let $\Delta$ now be the Hodge Laplacian of a compact Riemannian manifold, let
$\Pi_{\le\Lambda}$ be its spectral projector onto eigenvalues at most $\Lambda$, and let
$V_\Lambda=\ran\Pi_{\le\Lambda}$, a finite-dimensional space. For a family of operators
$S^{\kappa}_{e_i,e_j}$ obtained from a kernel by fixing two of its arguments, define
\[
 \Energy_d^{\Lambda}(\kappa)
 =\sum_{i,j}\Bigl(
 \norm{[d,S^{\kappa}_{e_i,e_j}]\Pi_{\le\Lambda}}_{\mathrm{HS}}^{2}
 +\norm{[d^{*},S^{\kappa}_{e_i,e_j}]\Pi_{\le\Lambda}}_{\mathrm{HS}}^{2}\Bigr).
\]

\begin{proposition}[Truncated compatibility]
\label{prop:truncated}
If $\Energy_d^{\Lambda}(\kappa)=0$ then every $S^{\kappa}_{e_i,e_j}$ commutes with $d$ and
$d^{*}$ \emph{on $V_\Lambda$}, and therefore descends to the cohomology of the truncated
complex $\Pi_{\le\Lambda}C^{\bullet}$ by Proposition~\ref{prop:descent}.
\end{proposition}

\begin{remark}[The truncation is not the continuum]
\label{rem:truncation}
Proposition~\ref{prop:truncated} is a statement about $V_\Lambda$ only. It gives no
information about $[d,S^{\kappa}]$ on the orthogonal complement of $V_\Lambda$, and
therefore none about the untruncated complex. Vanishing of $\Energy_d^{\Lambda}$ for every
$\Lambda$ would be a different and much stronger hypothesis. Furthermore, a numerically
located local minimum of $\Energy_d^{\Lambda}$ does not imply
$\Energy_d^{\Lambda}=0$; that conclusion requires an additional surjectivity or
controllability hypothesis on the variations of the defect map, which we do not have.
\end{remark}

\section{An induced Markov operator: a conditional construction}
\label{sec:markov}

This section is a construction under a hypothesis that is \emph{not verified here for any
kernel}. It is included because it is the shape a diffusion-theoretic extension would
take, and because stating it precisely makes clear what would have to be proved.

\begin{construction}[Induced Markov operator]
\label{constr:markov}
Suppose that a declared contraction of a kernel produces a measurable
$\mathcal W:X\times X\to[0,\infty)$ with $\mathcal W(p,s)=\mathcal W(s,p)$, and set
\[
 \mathsf d(p)=\int_X\mathcal W(p,s)\,d\nu(s).
\]
Assume $0<\mathsf d(p)<\infty$ for $\nu$-almost every $p$. Define
$\mathcal P(p,s)=\mathcal W(p,s)/\mathsf d(p)$, the measure $d\lambda=\mathsf d\,d\nu$, the
operator $(\mathcal Pf)(p)=\int_X\mathcal P(p,s)f(s)\,d\nu(s)$, and
$\Delta_{\mathcal P}=I-\mathcal P$.
\end{construction}

\begin{proposition}[Self-adjointness and positivity]
\label{prop:markov}
Under the hypotheses of Construction~\ref{constr:markov}, $\mathcal P$ is self-adjoint on
$L^{2}(X,\lambda)$, and for $f\in L^{2}(X,\lambda)$
\[
 \langle f,\Delta_{\mathcal P}f\rangle_{L^{2}(\lambda)}
 =\tfrac12\iint_{X\times X}\lvert f(p)-f(s)\rvert^{2}\,\mathcal W(p,s)\,d\nu(p)\,d\nu(s)
 \;\ge\;0 .
\]
\end{proposition}

\begin{proof}
For $f,g\in L^{2}(X,\lambda)$,
\[
 \langle\mathcal Pf,g\rangle_{L^{2}(\lambda)}
 =\int_X\Bigl(\int_X\frac{\mathcal W(p,s)}{\mathsf d(p)}f(s)\,d\nu(s)\Bigr)
 \overline{g(p)}\,\mathsf d(p)\,d\nu(p)
 =\iint \mathcal W(p,s)f(s)\overline{g(p)}\,d\nu(s)\,d\nu(p),
\]
which is symmetric in the pair $(f,g)$ after exchanging $p\leftrightarrow s$ and using
$\mathcal W(p,s)=\mathcal W(s,p)$; hence $\mathcal P$ is self-adjoint. Expanding,
\[
 \tfrac12\iint\lvert f(p)-f(s)\rvert^{2}\mathcal W(p,s)
 =\int_X\lvert f(p)\rvert^{2}\mathsf d(p)\,d\nu(p)
 -\operatorname{Re}\iint\mathcal W(p,s)f(s)\overline{f(p)},
\]
using symmetry of $\mathcal W$ for both cross terms, and the right-hand side is
$\langle f,f\rangle_{L^{2}(\lambda)}-\langle\mathcal Pf,f\rangle_{L^{2}(\lambda)}$, which
is $\langle f,\Delta_{\mathcal P}f\rangle_{L^{2}(\lambda)}$.
\end{proof}

The operator $\Delta_{\mathcal P}$ is the graph Laplacian of the weight $\mathcal W$; the
Dirichlet form above is the standard one, and the diffusion-map construction
\cite{CoifmanLafon2006} is the analytic setting in which such operators are usually
studied.

\begin{remark}[Spectral dimension is a further hypothesis]
\label{rem:spectraldim}
If $e^{-t\Delta_{\mathcal P}}$ is trace class one may set $\Theta(t)=\Tr
e^{-t\Delta_{\mathcal P}}$, and \emph{if} $\Theta(t)\sim Ct^{-d_s/2}$ as $t\downarrow0$ one
may call $d_s=-2\lim_{t\downarrow0}\frac{d\log\Theta(t)}{d\log t}$ the spectral dimension.
Both the trace-class property and the existence of the limit are additional hypotheses.
Neither is a consequence of a kernel being square-integrable.
\end{remark}

\begin{remark}[The hypothesis of Construction~\ref{constr:markov} is unverified]
\label{rem:markovunverified}
We exhibit no contraction of a kernel $\kappa\in L^{2}(X^{4})$ that is guaranteed to
produce a symmetric nonnegative $\mathcal W$ with $0<\mathsf d<\infty$ almost everywhere,
and we give no condition on $\kappa$ under which one exists. Until such a condition is
supplied, Construction~\ref{constr:markov} has no verified instance.
\end{remark}

\section{Multiscale transport defects}
\label{sec:multiscale}

Let $(V_N,\mu_N,P_N)$ be a family indexed by a resolution parameter $N$, with prolongation
maps $J_N^{M}:V_N\to V_M$ and restriction maps $R_M^{N}:V_M\to V_N$. Define the
\emph{transport defects}
\[
 \mathcal D_\mu^{N,M}=J_N^{M}\mu_N-\mu_M(J_N^{M}\cdot,\ldots,J_N^{M}\cdot),
 \qquad
 \mathcal D_P^{N,M}=P_MJ_N^{M}-J_N^{M}P_N,
\]
and compare subspaces by $\norm{P_M-J_N^{M}P_NR_M^{N}}$ or by principal angles.

These are definitions. No theorem is asserted about them here.

\begin{remark}[A decreasing sequence is not a limit theorem]
\label{rem:nolimit}
A finite decreasing sequence of transport defects does not establish a continuum limit.
Establishing one would additionally require: explicit topologies on the family; uniform
bounds on $\norm{\mu_N}_{\op}$, on $\norm{P_N}$ and on the leakage norms $\rho_N$; a
compactness argument; convergence of the maps; convergence of the projectors in a stated
operator topology; stability of whichever identities are of interest; and identification
of the limiting operator. See Question~\ref{q:continuum}.
\end{remark}

\section{Open analytical questions}
\label{sec:open}

The following are questions. None of them is a result of this article, and nothing above
depends on any of them.

\begin{question}[Continuum limit]
\label{q:continuum}
Under what hypotheses on a family $(V_N,\mu_N,P_N)$ with prolongations $J_N^{M}$ do the
unprojected and recursively projected evaluations of a fixed finite tree converge, and does
the estimate $E_T^{\proj}\le(k-1)\rho M^{k-1}L_T$ of \eqref{eq:companion} pass to the
limit? A plausible set of hypotheses is: $J_N^{M}\mu_N(R\,\cdot,\ldots,R\,\cdot)\to\mu$ and
$J_N P_N R_N\to P$ in a stated topology, together with
$\sup_N\norm{\mu_N}_{\op}<\infty$, $\sup_N\norm{P_N}\le1$ and $\sup_N\rho_N<\infty$. For
trees whose size grows with $N$ a further control on depth and on the summability of the
sources would be needed.
\end{question}

\begin{question}[Membership in a multilinear pseudodifferential class]
\label{q:psido}
Under what hypotheses does a kernel-defined operator $\Kop_\kappa$ admit a local symbol
$a(x,\xi_1,\ldots,\xi_a)$ satisfying estimates of the form
\[
 \bigl\lvert\partial_x^{\alpha}\partial_{\xi_1}^{\beta_1}\cdots
 \partial_{\xi_a}^{\beta_a}a\bigr\rvert
 \le C_{\alpha,\beta}\Bigl(1+\textstyle\sum_j\lvert\xi_j\rvert\Bigr)^{m-\varrho\lvert\beta\rvert+\delta\lvert\alpha\rvert},
\]
so that it belongs to a multilinear H\"ormander class? Square-integrability of $\kappa$ is
far from sufficient. The multilinear Calder\'on--Zygmund theory of Grafakos and Torres
\cite{GrafakosTorres2002} and the bilinear H\"ormander classes of B\'enyi, Maldonado,
Naibo and Torres \cite{BenyiMaldonadoNaiboTorres2010} indicate what such a theory requires:
symbol classes, invariance under changes of coordinates, Sobolev continuity, composition,
adjoints, and stability of the class under the formation of associators. A realistic first
target would be a compact manifold, kernels smooth off the diagonal, symbols of order zero
and spectral projectors as the $P_\tau$.
\end{question}

\begin{question}[Microlocal regularity]
\label{q:microlocal}
If Question~\ref{q:psido} were answered affirmatively, what containment of the form
$\mathrm{WF}\bigl(\Kop_\kappa(f_1,\ldots,f_a)\bigr)\subseteq
\mathcal R_\kappa\bigl(\mathrm{WF}(f_1),\ldots,\mathrm{WF}(f_a)\bigr)$ holds, for which
canonical relation $\mathcal R_\kappa$? Does orthogonal projection create singularities;
is the orthogonal component of the error more regular than the projected one; and is there
a microlocal analogue of the passage from $k$ to $k-1$ sources?
\end{question}

\begin{question}[Effect of projection on the class]
\label{q:projclass}
Is a class of kernel-defined operators stable under the operation
$\mathcal S\mapsto P\mathcal S(P\,\cdot,\ldots,P\,\cdot)$ that the finite theory performs
at every vertex, and how does the leakage norm $\rho$ behave under it?
\end{question}

\begin{question}[Verifying Construction~\ref{constr:markov}]
\label{q:markov}
Give a condition on $\kappa\in L^{2}(X^{4})$ under which some declared contraction yields a
symmetric nonnegative weight with $0<\mathsf d<\infty$ almost everywhere.
\end{question}

We note explicitly that holonomic $D$-modules, the Riemann--Hilbert correspondence and
algebraisation of a limiting projector have sometimes been mentioned in connection with
this circle of ideas. No hypothesis, no statement and no proof relating them to the
material above exists in this article, and they are therefore not listed even as
questions.

\section{Summary of status}
\label{sec:status}

\begin{table}[h]
\centering\small
\caption{Status of each statement in this article.}
\label{tab:status}
\begin{tabularx}{\linewidth}{@{}lX@{}}
\toprule
Statement & Status \\
\midrule
Prop.~\ref{prop:bounded}, \ref{prop:stability} & proved \\
Lemma~\ref{lem:composite} & proved; closes a gap in an earlier formulation \\
Prop.~\ref{prop:represents} & proved, using Lemma~\ref{lem:composite} \\
Prop.~\ref{prop:converse} & proved under separability of $L^{2}(X,\nu)$ \\
Thm.~\ref{thm:commutator} & proved; an identity for any bilinear operation \\
Def.~\ref{def:ralg} & a definition, not a theorem; see Remark~\ref{rem:notcurvature} \\
Prop.~\ref{prop:descent}, \ref{prop:hodge} & proved for finite complexes; standard \\
Prop.~\ref{prop:truncated} & proved for the truncation only; see Remark~\ref{rem:truncation} \\
Constr.~\ref{constr:markov}, Prop.~\ref{prop:markov} & conditional construction; hypothesis unverified \\
Remark~\ref{rem:spectraldim} & definition under two further hypotheses \\
Section~\ref{sec:multiscale} & definitions only \\
Section~\ref{sec:open} & open questions \\
Eq.~\eqref{eq:companion} & proved in \cite{ChaveroJassoTrees}; not reproved here \\
\bottomrule
\end{tabularx}
\end{table}

None of the statements above has been checked by anyone other than the author, and no
assessment of originality has been carried out for any of them. Several --- notably
Propositions~\ref{prop:descent} and~\ref{prop:hodge}, and the Stiefel gradient of
Section~\ref{sec:variational} --- are standard, and are included for completeness with
citations rather than as contributions.

\bibliographystyle{amsplain}
\bibliography{references}

\end{document}
